\documentclass[11pt,a4paper]{article}

\usepackage[T1]{fontenc}
\usepackage[utf8]{inputenc}
\usepackage{lmodern}
\usepackage[margin=2.6cm]{geometry}
\usepackage{amsmath,amssymb,amsthm,mathtools}
\usepackage{microtype}
\usepackage[hidelinks]{hyperref}

\newtheorem{theorem}{Theorem}
\newtheorem{lemma}{Lemma}

\newcommand{\HH}{\mathcal H_k}
\newcommand{\E}{\mathbb E}
\newcommand{\Prob}{\mathbb P}
\newcommand{\norm}[1]{\left\lVert #1\right\rVert}
\newcommand{\inner}[2]{\left\langle #1,#2\right\rangle}
\newcommand{\op}{\mathrm{op}}
\newcommand{\HS}{\mathrm{HS}}

\title{Pairwise Guarantees for Directed\\
Regularized Pearson-Type Classification}
\author{}
\date{}

\begin{document}
\maketitle

\section{Setting and score}

Let $\HH$ be a real, separable RKHS with measurable
feature map $\phi(z)=k(\cdot,z)$. We work with
positive-semidefinite, pointwise nonnegative kernels
satisfying
\begin{equation}\label{eq:kernel}
    k(x,y)\ge0,
    \qquad
    \sup_z k(z,z)\le K<\infty,
    \qquad K>0.
\end{equation}
Nonnegativity implies
\begin{equation}\label{eq:feature-difference}
    \norm{\phi(x)-\phi(y)}_{\HH}^2
    =k(x,x)+k(y,y)-2k(x,y)\le2K,
\end{equation}
which determines the concentration constants below.

Fix a reference distribution $P_i$ and an incoming
distribution $Q$. Let
\[
    Z_{i,1},\ldots,Z_{i,n}
    \overset{\mathrm{i.i.d.}}{\sim}P_i,
    \qquad
    Z'_1,\ldots,Z'_n
    \overset{\mathrm{i.i.d.}}{\sim}Q,
\]
with the two collections independent. Write
\[
    X_i=(Z_{i,1},\ldots,Z_{i,n}),
    \qquad
    X'=(Z'_1,\ldots,Z'_n).
\]

Define
\begin{align*}
    \mu_P&=\E_{Z\sim P}\phi(Z),
    &
    T_P&=\E_{Z\sim P}[\phi(Z)\otimes\phi(Z)],\\
    \widehat\mu_i&=\frac1n\sum_{r=1}^n\phi(Z_{i,r}),
    &
    \widehat\mu'&=\frac1n\sum_{r=1}^n\phi(Z'_r),\\
    \widehat T_i&=\frac1n\sum_{r=1}^n
        \phi(Z_{i,r})\otimes\phi(Z_{i,r}),
    &
    \widehat\Delta_i&=\widehat\mu'-\widehat\mu_i.
\end{align*}
Here $(u\otimes v)w=\inner{v}{w}_{\HH}u$,
$T_P$ is the uncentered second-moment operator,
and $\Delta_i=\mu_Q-\mu_{P_i}$.

For fixed $\lambda>0$, consider
\begin{align}
    \chi^2_\lambda(Q;P_i)
    &:=
    \frac12\inner{\Delta_i}
    {(T_{P_i}+\lambda I)^{-1}\Delta_i}_{\HH},
    \label{eq:population}\\
    \widehat\chi^2_{\lambda,i}
    &:=
    \frac12\inner{\widehat\Delta_i}
    {(\widehat T_i+\lambda I)^{-1}
    \widehat\Delta_i}_{\HH}.
    \label{eq:empirical}
\end{align}
This is a directed regularized RKHS Pearson-type
discrepancy, with the reference distribution after
the semicolon. It need not equal the classical
Pearson divergence.

The kernel is fixed before sampling. A separate fixed
kernel may be assigned to each reference, provided
both windows in a comparison are evaluated using
that reference's kernel. All guarantees below concern
one fixed comparison and are over both sampled windows.

\subsection{Concentration}

For $\gamma\in(0,1)$, set
\begin{equation}\label{eq:radii}
    b_n(\gamma)
    :=
    \sqrt{\frac{2K}{n}}
    \left(1+\sqrt{\log\frac1\gamma}\right),
    \qquad
    e_n(\gamma)
    :=
    \frac K{\sqrt n}
    \left(1+\sqrt{\log\frac1\gamma}\right).
\end{equation}

\begin{lemma}\label{lem:concentration}
For every $\gamma\in(0,1)$,
\begin{align}
    \Prob\left\{
        \norm{\widehat\Delta_i-\Delta_i}_{\HH}
        \le b_n(\gamma)
    \right\}
    &\ge1-\gamma,
    \label{eq:mean-concentration}\\
    \Prob\left\{
        \norm{\widehat T_i-T_{P_i}}_{\op}
        \le e_n(\gamma)
    \right\}
    &\ge1-\gamma.
    \label{eq:operator-concentration}
\end{align}
\end{lemma}

\begin{proof}
Independence and centering give
\begin{align*}
    \E\norm{\widehat\Delta_i-\Delta_i}_{\HH}^2
    &=
    \frac1n
    \left(
        \E_{Z\sim Q}k(Z,Z)-\norm{\mu_Q}_{\HH}^2
    \right)\\
    &\quad+
    \frac1n
    \left(
        \E_{Z\sim P_i}k(Z,Z)
        -\norm{\mu_{P_i}}_{\HH}^2
    \right)\\
    &\le\frac{2K}{n}.
\end{align*}
Thus the expected error norm is at most
$\sqrt{2K/n}$.
By \eqref{eq:feature-difference}, replacing one of
the $2n$ observations changes that norm by at most
$\sqrt{2K}/n$. McDiarmid's inequality therefore yields,
for $t>0$,
\[
    \Prob\left\{
        \norm{\widehat\Delta_i-\Delta_i}_{\HH}
        -
        \E\norm{\widehat\Delta_i-\Delta_i}_{\HH}
        \ge t
    \right\}
    \le
    \exp\left(-\frac{nt^2}{2K}\right).
\]
Taking $t=\sqrt{(2K/n)\log(1/\gamma)}$ proves
\eqref{eq:mean-concentration}.

For the operator, write $V(z)=\phi(z)\otimes\phi(z)$.
Since $\norm{V(z)}_{\HS}=k(z,z)\le K$,
\[
    \E\norm{\widehat T_i-T_{P_i}}_{\HS}^2
    =
    \frac{
        \E_{Z\sim P_i}\norm{V(Z)}_{\HS}^2
        -
        \norm{T_{P_i}}_{\HS}^2
    }{n}
    \le\frac{K^2}{n}.
\]
Moreover,
\[
    \norm{V(x)-V(y)}_{\HS}^2
    =
    k(x,x)^2+k(y,y)^2-2|k(x,y)|^2
    \le2K^2.
\]
Consequently, the norm
$G=\norm{\widehat T_i-T_{P_i}}_{\HS}$
has expectation at most $K/\sqrt n$ and replacement
sensitivity at most $\sqrt2K/n$ for each of its $n$
observations. McDiarmid's inequality gives
\[
    \Prob\{G-\E G\ge t\}
    \le
    \exp\left(-\frac{nt^2}{K^2}\right).
\]
Taking $t=K\sqrt{\log(1/\gamma)/n}$ and using
$\norm{\cdot}_{\op}\le\norm{\cdot}_{\HS}$
proves \eqref{eq:operator-concentration}.
\end{proof}

\section{Upper bound under the null hypothesis}

\begin{theorem}[Null upper bound]\label{thm:null}
If $Q=P_i$, then for every $\rho_0\in(0,1)$,
\[
    \Prob\left\{
        0\le\widehat\chi^2_{\lambda,i}
        \le U_n(\rho_0)
    \right\}
    \ge1-\rho_0,
\]
where
\begin{equation}\label{eq:U}
    \boxed{
    U_n(\rho_0)
    :=
    \frac{b_n(\rho_0)^2}{2\lambda}
    =
    \frac{K}{n\lambda}
    \left(1+\sqrt{\log\frac1{\rho_0}}\right)^2.
    }
\end{equation}
\end{theorem}

\begin{proof}
Since $\widehat T_i\succeq0$,
\[
    0\le\widehat\chi^2_{\lambda,i}
    \le\frac1{2\lambda}
        \norm{\widehat\Delta_i}_{\HH}^2.
\]
Under the null, $\Delta_i=0$.
Applying \eqref{eq:mean-concentration} at level
$\rho_0$ proves the result.
\end{proof}

\section{Lower bound under a separated alternative}

\subsection{Separation assumption}

For a fixed pair, assume
\begin{equation}\label{eq:separation}
    \chi^2_\lambda(Q;P_i)\ge\delta>0.
\end{equation}
Alternatively, for a family $\mathcal P$ of admissible
regimes, one may assume
\[
    \delta_{\min}
    :=
    \inf_{\substack{P,Q\in\mathcal P\\P\ne Q}}
        \chi^2_\lambda(Q;P)>0
\]
and take $\delta=\delta_{\min}$.
Each score uses its reference kernel when kernels
are reference-specific. Any novel distribution
covered by this uniform guarantee must belong
to $\mathcal P$.

\begin{theorem}[Alternative lower bound]
\label{thm:alternative}
Under \eqref{eq:separation}, for every
$\rho_1\in(0,1)$,
\[
    \Prob\left\{
        \widehat\chi^2_{\lambda,i}
        \ge L_n(\delta,\rho_1)
    \right\}
    \ge1-\rho_1,
\]
where
\begin{equation}\label{eq:L}
    \boxed{
    L_n(\delta,\rho_1)
    :=
    \frac{
        \left[
            \sqrt{2\delta}
            -b_n(\rho_1/2)/\sqrt\lambda
        \right]_+^2
    }{
        2\left(1+e_n(\rho_1/2)/\lambda\right)
    }.
    }
\end{equation}
Here $[x]_+=\max\{x,0\}$.
\end{theorem}

\begin{proof}
Set $b=b_n(\rho_1/2)$ and $e=e_n(\rho_1/2)$.
By Lemma~\ref{lem:concentration} and the union bound,
with probability at least $1-\rho_1$,
\[
    \norm{\widehat\Delta_i-\Delta_i}_{\HH}\le b,
    \qquad
    \norm{\widehat T_i-T_{P_i}}_{\op}\le e.
\]
Work on this event and write
\[
    A=T_{P_i}+\lambda I,
    \qquad
    \widehat A=\widehat T_i+\lambda I,
    \qquad
    c=1+\frac e\lambda.
\]
Since $A\succeq\lambda I$,
\[
    \widehat A\preceq A+eI\preceq cA,
    \qquad
    \widehat A^{-1}\succeq c^{-1}A^{-1}.
\]
Using the norm
$\norm{v}_{A^{-1}}^2=\inner{v}{A^{-1}v}_{\HH}$,
the reverse triangle inequality gives
\begin{align*}
    \sqrt{2\widehat\chi^2_{\lambda,i}}
    &\ge
    \frac1{\sqrt c}\norm{\widehat\Delta_i}_{A^{-1}}\\
    &\ge
    \frac1{\sqrt c}
    \left[
        \norm{\Delta_i}_{A^{-1}}
        -
        \norm{\widehat\Delta_i-\Delta_i}_{A^{-1}}
    \right]_+\\
    &\ge
    \frac1{\sqrt c}
    \left[
        \sqrt{2\delta}-\frac b{\sqrt\lambda}
    \right]_+.
\end{align*}
The last inequality uses \eqref{eq:separation} and
$A^{-1}\preceq\lambda^{-1}I$.
Squaring and dividing by $2$ proves \eqref{eq:L}.
\end{proof}

\subsection{Threshold interval}

Accept the reference if
$\widehat\chi^2_{\lambda,i}\le\eta$
and reject it otherwise.
Theorems~\ref{thm:null}--\ref{thm:alternative}
give the respective error bounds
\[
    \Prob_{H_0}
    \{\widehat\chi^2_{\lambda,i}>\eta\}\le\rho_0,
    \qquad
    \Prob_{H_1}
    \{\widehat\chi^2_{\lambda,i}\le\eta\}\le\rho_1,
\]
whenever
\begin{equation}\label{eq:threshold-interval}
    U_n(\rho_0)\le\eta<L_n(\delta,\rho_1).
\end{equation}
Here $H_0$ denotes $Q=P_i$, and $H_1$ denotes any
fixed pair satisfying \eqref{eq:separation}.

Writing
\[
    b_0=b_n(\rho_0),
    \qquad
    b_1=b_n(\rho_1/2),
    \qquad
    c=1+\frac{e_n(\rho_1/2)}{\lambda},
\]
the interval is nonempty if and only if
\begin{equation}\label{eq:threshold-separation}
    \boxed{
    \delta>\frac{(b_1+\sqrt c\,b_0)^2}{2\lambda}.
    }
\end{equation}
Indeed, since $b_0>0$, the inequality
$L_n(\delta,\rho_1)>U_n(\rho_0)$ is equivalent to
\[
    \sqrt{2\delta}
    >
    \frac{b_1+\sqrt c\,b_0}{\sqrt\lambda}.
\]

\subsection{Threshold calibration and required separation}

Choose a common error tolerance
$\rho_0=\rho_1=\rho\in(0,1)$ and set
$\eta=U_n(\rho)$. In terms of the initial parameters,
\begin{equation}\label{eq:explicit-threshold}
    \boxed{
    \eta
    =
    \frac{K}{n\lambda}
    \left(1+\sqrt{\log\frac1\rho}\right)^2.
    }
\end{equation}
This guarantees a null rejection probability at most
$\rho$, independently of the separation assumption.

Substitution into \eqref{eq:threshold-separation}
gives the separation required for a nonempty
threshold interval:
\begin{equation}\label{eq:explicit-separation}
\boxed{
\begin{aligned}
    \delta>
    \frac{K}{n\lambda}
    \Bigg[
        &1+\sqrt{\log\frac2\rho}\\
        &+
        \left(1+\sqrt{\log\frac1\rho}\right)
        \sqrt{
            1+\frac{K}{\lambda\sqrt n}
            \left(1+\sqrt{\log\frac2\rho}\right)
        }
    \Bigg]^2.
\end{aligned}
}
\end{equation}
Under this condition, every fixed alternative pair
with $\chi^2_\lambda(Q;P_i)\ge\delta$ satisfies
\[
    \Prob\left\{
        \widehat\chi^2_{\lambda,i}\le\eta
    \right\}\le\rho.
\]

Condition~\eqref{eq:explicit-separation} is equivalent
to nonemptiness of the interval defined by our bounds.
It is sufficient, but not necessary, for successful
classification. These guarantees concern one fixed
reference--incoming comparison; simultaneous control
over the memory bank and repeated online decisions
requires a separate analysis.

\end{document}